\chapter{开始之前}
本书面向会运行 Python 脚本、认识 HTML 标签和属性的读者。第一次接触 \texttt{venv} 也不要紧：本节把环境、源码和第一条结果接起来。示例只访问本地文件或本机回环地址，运行时不需要登录任何网站。

\section*{随书文件在哪里}
本 PDF 配有同版随书源码包，GitHub Release 中的文件名为 \nolinkurl{Python-Web-Scraping-Source.zip}；解压后进入 \texttt{book\_crawler} 目录。若只收到 PDF，请同时取得该源码包。源码包中保留了正文所展示的独立文件，PDF 和代码以同一版次为准。

\begin{Verbatim}[numbers=none]
book_crawler/
  main.tex, preamble.tex       # 讲义入口与版式
  examples/ch01_links.py       # 第一章示例
  examples/common.py           # 本地 HTTP 服务的启动与关闭
  examples/validation.py       # 响应及分页契约的可测函数
  examples/fixtures/           # Scrapy 使用的固定页面
  tests/                       # 成功和失败路径的自动化断言
  requirements-lock.txt        # Python 3.12 完整依赖版本
  examples/package-lock.json   # AST 示例的 Node 依赖版本
\end{Verbatim}

代码文件与 PDF 的例号同章对应，附录 B 列出路径和期望结果。\texttt{build/} 保存本地运行结果，不需要手工创建；程序会在需要时建立。执行命令前先确认终端当前目录含 \texttt{main.tex}，否则相对路径会指错地方。

\section*{安装与第一条结果}
已验证环境为 Python 3.12、Node.js 22、XeLaTeX 和项目锁文件所列的依赖。只读第一章时，标准库就够了；完成全书示例再安装依赖。Windows PowerShell：
\begin{Verbatim}[numbers=none]
cd book_crawler
py -3.12 -m venv .venv
.\.venv\Scripts\python.exe -m pip install -r requirements-lock.txt
.\.venv\Scripts\python.exe examples/ch01_links.py
\end{Verbatim}
最后一行应得到两条链接，顺序如下；第二次出现的同一文章链接会被保序去重。
\begin{Verbatim}[numbers=none]
https://example.org/articles/intro
https://example.org/articles/advanced
\end{Verbatim}
Linux 或 macOS 将 Python 可执行文件换为 \texttt{.venv/bin/python}；首次建立虚拟环境可用 \texttt{python3.12 -m venv .venv}。上面的 \texttt{example.org} 只是固定 HTML 中的链接文字，第一章不会向它发请求。

\section*{浏览器、Node 与回归测试}
第 10 章需要 Playwright 浏览器。Windows 已安装 Edge 时，在同一 PowerShell 窗口指定浏览器；否则安装 Chromium 并清除该环境变量：
\begin{Verbatim}[numbers=none]
$env:BROWSER_CHANNEL = 'msedge'
# 若不用 Edge，则执行：
Remove-Item Env:BROWSER_CHANNEL
.\.venv\Scripts\python.exe -m playwright install chromium
\end{Verbatim}
第 19 章需 Node.js 22。安装并运行自写的 AST 示例：
\begin{Verbatim}[numbers=none]
cd examples
npm ci
cd ..
node examples/ch19_ast.js
\end{Verbatim}

\begin{Verbatim}[numbers=none]
.\.venv\Scripts\python.exe -m pytest -q
\end{Verbatim}
测试既检查预期输出，也检查空字段、404、重复游标和中断恢复等失败样本。常见启动问题按下面顺序排查：
\begin{itemize}
\item 缺少模块：核对 Python 是否来自新虚拟环境，再安装锁文件。
\item 找不到 \texttt{examples/}：回到含 \texttt{main.tex} 的项目根目录。
\item 浏览器启动失败：按上一段安装 Chromium，或指定本机 Edge。
\end{itemize}
依赖缺失不是“网站反爬”；这些示例根本没去外站。
